\documentclass[10pt]{amsart}
\usepackage{calc,amssymb,amsmath,ulem,stmaryrd,fullpage}
\usepackage{enumerate}
\usepackage{alltt}
\usepackage{multicol}
\RequirePackage[dvipsnames,usenames]{color}
\let\mffont=\sf
\normalem
%\input{kmacros3.sty}
\input{xy}
\xyoption{all}
\usepackage{hyperref}
\usepackage{graphicx}
\usepackage[all,cmtip]{xy}
\usepackage{verbatim}
\usepackage[left=0.95in,top=0.95in,right=0.95in,bottom=0.95in]{geometry}
\newcommand{\tauCohomology}{T}
\newcommand{\FCohomology}{S}


\theoremstyle{theorem}
%\newtheorem*{mainthm*}{Main Theorem}

\newcommand{\note}[1]{\marginpar{\sffamily\tiny #1}}

\def\todo#1{\textcolor{Mahogany}%
{\sffamily\footnotesize\newline{\color{Mahogany}\fbox{\parbox{\textwidth-15pt}{#1}}}\newline}}

\renewcommand{\O}{\mathcal O}
\newcommand{\ie}{{\itshape i.e.} }
\newcommand{\roundup}[1]{\lceil #1 \rceil}
\newcommand{\rounddown}[1]{\lfloor #1 \rfloor}
\renewcommand{\to}[1][]{\xrightarrow{\ #1\ }}
\newcommand{\onto}[1][]{\protect{\xrightarrow{\ #1\ }\hspace{-0.8em}\rightarrow}}
\newcommand{\into}[1][]{\lhook \joinrel \xrightarrow{\ #1\ }}
%\newcommand{DBDual}[1]{{\underline{\omega}_{#1}}}
\newcommand{\DBDual}[1]{{{\uuline \omega}{}^{\mydot}_{#1}}}
\newcommand{\Tt}{{\mathfrak{T}}}
%\newcommand{\bn}{{\mathfrak{n}} }
\newcommand{\sol}[1]{ \vskip 6pt{\color{blue} {\bf Solution:} #1} \vskip 6pt}

\newcommand{\bR}{{\mathbb{R}}}
\newcommand{\va}{{\vec{a}}}
\newcommand{\vb}{{\vec{b}}}
\newcommand{\vzero}{{\vec{0}}}
\newcommand{\vi}{{\vec{i}}}
\newcommand{\vj}{{\vec{j}}}
\newcommand{\vk}{{\vec{k}}}
\newcommand{\vh}{{\vec{h}}}
\newcommand{\vr}{{\vec{r}}}
\newcommand{\vu}{{\vec{u}}}
\newcommand{\vv}{{\vec{v}}}
\newcommand{\vw}{{\vec{w}}}
\newcommand{\vx}{{\vec{x}}}
\newcommand{\vy}{{\vec{y}}}
\newcommand{\vz}{{\vec{z}}}
\newcommand{\vT}{{\vec{T}}}
\newcommand{\vN}{{\vec{N}}}
\newcommand{\vF}{{\vec{F}}}
\newcommand{\vG}{{\vec{G}}}
\newcommand{\bF}{\mathbb{F}}
\newcommand{\bQ}{\mathbb{Q}}
\newcommand{\bZ}{\mathbb{Z}}
\newcommand{\bC}{\mathbb{C}}


\begin{document}
\title{Homework \#3 -- Math 5320\\Spring 2020}
\author{Due Wednesday, March 18th, in Gradescope}

\maketitle


\noindent
{\bf 1.}  Let $\alpha$ be a complex root of the polynomial $x^3 - 3x + 4$.  Write:
\[
    1 \over \alpha^2 + \alpha + 1
\]
as a linear combination of $1, \alpha, \alpha^2$.  
\newpage
\noindent
{\bf 2.}  Let $a$ be a positive rational number that is not a square in $\bQ$.  Prove that $a^{1/4}$ has degree $4$ over $\bQ$.
\newpage
\noindent
{\bf 3.}  Let $\alpha, \beta \in \bC$ be roots of irreducible polynomials $f(x), g(x) \in \bQ[x]$ respectively.  Let $K = \bQ[\alpha]$ and $L = \bQ[\beta]$.  Prove that $f(x)$ is irreducible in $L[x]$ if and only if $g(x)$ is irreducible in $K[x]$.  
\newpage
\noindent
{\bf 4.}  Determine the irreducible polynomial for $\alpha = \sqrt{3} + \sqrt{5}$ over the following fields.
{\bf (a)}  $\bQ$
\vskip 3pt
{\bf (b)}  $\bQ[\sqrt{5}]$
\vskip 3pt
{\bf (c)}  $\bQ[\sqrt{10}]$
\vskip 3pt
{\bf (d)}  $\bQ[\sqrt{15}]$
\newpage
\noindent
{\bf 5.}  Is $i$ in the field $\bQ[(-2)^{1/4}]$?
\newpage
\noindent
{\bf 6.}  Suppose that $F$ is a field.  Prove that if $[F[\alpha] : F]$ is finite and odd, then $F[\alpha^2] = F[\alpha]$.
\newpage
\noindent
{\bf 7.}  Prove that $\bQ[\sqrt{2} + \sqrt{3}] = \bQ[\sqrt{2}, \sqrt{3}]$.
\newpage
\noindent
{\bf 8.}  Suppose $\alpha_1, \dots, \alpha_n \in \bC$ are such that $\alpha_i^2 \in \bQ$ for all $i$.  Prove that $2^{1/3}$ is not in $\bQ[\alpha_1, \dots, \alpha_n]$.
\newpage
\noindent
{\bf 9.}  Suppose that $f(x)$ is an irreducible polynomial of degree $n$ over a field $F$.  Let $g(x)$ be any polynomial in $F[x]$.  Prove that every irreducible factor of the composite polynomial $f(g(x))$ has degree divisible by $n$.  



\end{document} 